% 18-comptime.tex — Chapter 18: Compile-Time Evaluation and Generics
\chapter{Compile-Time Evaluation and Generics}
\label{chap:18-comptime}
\index{comptime}
\index{compile-time evaluation}
\index{generics}

\pnum
Lain's compile-time evaluation (CTFE) system evaluates a subset of the AST
during compilation to support generic functions, type aliases, and comptime-if
branches.  The CTFE engine is in \srcref{src/sema/comptime.h}; generic
instantiation is in \srcref{src/sema/generic.h}.

% -------------------------------------------------------------------------
\section{Comptime environment}
\label{sec:ct-env}
\index{ComptimeEnv}
% -------------------------------------------------------------------------

\pnum
The CTFE interpreter maintains a linked list of bindings:

\begin{ccode}
typedef struct ComptimeEnv {
    Id       *name;
    Expr     *value;
    struct ComptimeEnv *next;
} ComptimeEnv;
\end{ccode}

\pnum
\cname{comptime\_env\_push} prepends a new binding to the front of the list.
\cname{comptime\_env\_lookup} performs a linear search by name.  The environment
is passed by value through the recursive evaluator; there is no heap-allocated
scope stack.

% -------------------------------------------------------------------------
\section{Expression evaluation}
\label{sec:ct-eval}
\index{comptime\_evaluate\_expr}
% -------------------------------------------------------------------------

\pnum
\cname{comptime\_evaluate\_expr(arena, expr, env)} reduces an expression to a
value.  It handles:

\begin{itemize}
\item \cname{EXPR\_LITERAL} and \cname{EXPR\_STRING}: returned as-is.
\item \cname{EXPR\_IDENTIFIER}: looked up in \texttt{env}, then in globals.
  Struct/enum/type names are wrapped in an \cname{EXPR\_TYPE} node.
  Primitive type names (\lain{int}, \lain{u8}, etc.) are returned as canonical
  \cname{EXPR\_TYPE} nodes.
\item \cname{EXPR\_BINARY}: arithmetic operators are evaluated numerically and
  returned as \cname{EXPR\_LITERAL}.
\item \cname{EXPR\_CALL}: dispatched to \cname{comptime\_evaluate\_function}.
\end{itemize}

\pnum
Recursion depth is bounded by \cname{COMPTIME\_MAX\_DEPTH = 256} to prevent
stack overflow from runaway CTFE loops.  Exceeding this limit exits with
error \diagcode{E045}.

% -------------------------------------------------------------------------
\section{Generic function instantiation}
\label{sec:ct-generics}
\index{generic functions!instantiation}
% -------------------------------------------------------------------------

\pnum
A function with at least one \cname{comptime} parameter is a generic function.
When a call to such a function is resolved, the resolver:
\begin{enumerate}
\item Evaluates all comptime arguments at compile time.
\item Deep-clones the function body using \srcref{src/ast\_clone.h}.
\item Calls \cname{generic\_substitute\_expr} (from \srcref{src/sema/generic.h})
  to replace occurrences of the comptime parameter names with the evaluated
  types or values.
\item Registers the specialised clone as a new \cname{DECL\_FUNCTION} with a
  mangled name and inserts it into the global scope.
\item Resolves the specialised body recursively.
\end{enumerate}

\pnum
Because specialisation happens during name resolution, every generic function
appears in the emitted C output as one or more monomorphic functions with
mangled names derived from the argument types.  There is no runtime
polymorphism.

% -------------------------------------------------------------------------
\section{Comptime-if}
\label{sec:ct-if}
\index{comptime if}
% -------------------------------------------------------------------------

\pnum
A \cname{STMT\_COMPTIME\_IF} evaluates its condition at compile time via
\cname{comptime\_evaluate\_expr}.  The condition typically tests \cname{@os}
or \cname{@arch} builtin values against platform constants.  The resolver
records which branch was selected in \cname{StmtComptimeIf.is\_taken}; the
walk phase only analyses the selected branch.  The unselected branch is
never type-checked or emitted.

% -------------------------------------------------------------------------
\section{Type alias evaluation}
\label{sec:ct-typealiases}
\index{type aliases}
% -------------------------------------------------------------------------

\pnum
\lain{type Name = Expr} declarations are evaluated at resolve time.
\cname{comptime\_evaluate\_function} is called with the right-hand side
expression and the current environment.  For parameterised aliases
(e.g.\ \lain{type Option(comptime T) = ...}), the argument is substituted
before evaluation.  The result is stored in the \cname{DeclTypeAlias} node
and used by the emitter to generate the corresponding C \texttt{typedef}.
